\chapter{Delusions}

\section*{Definition}
Fixed, false beliefs that persist despite contradictory evidence, inconsistent with the individual's cultural or religious background, and representing a disturbance in thought content.

\section*{What Happens in Your Body}
Dysfunction in prefrontal cortex (reality testing and belief evaluation) and limbic system (salience attribution) leads to formation and maintenance of delusional beliefs. Dopamine hyperactivity in mesolimbic pathways contributes to aberrant salience where neutral stimuli acquire false meaning. Delusions are categorized by content: persecutory, grandiose, referential, somatic, erotomanic, and nihilistic.

\section*{Causes}
\begin{itemize}
  \item Schizophrenia or schizoaffective disorder
  \item Bipolar disorder (manic phase)
  \item Major depressive disorder with psychotic features
  \item Delirium (acute confusional state)
  \item Dementia (Alzheimer, Lewy body)
  \item Substance use or withdrawal (amphetamine, cocaine, alcohol withdrawal)
  \item Medication side effects (corticosteroids, anticholinergics, dopaminergic drugs)
  \item Metabolic or endocrine disorders (thyroid disease, B12 deficiency)
  \item Parkinson disease (medication-induced)
  \item Brain tumors or traumatic brain injury
\end{itemize}

\section*{When to See a Doctor}
See your doctor if:
\begin{itemize}
  \item Symptoms are severe or getting worse over time
  \item New-onset delusions, especially with agitation, self-harm risk, or inability to care for oneself (immediate psychiatric evaluation)
  \item You have other concerning symptoms such as fever, weight loss, or bleeding
\end{itemize}

\section*{Related Symptoms}
\begin{itemize}
  \item Hallucinations
  \item Agitation
  \item Confusion
  \item Suspiciousness
  \item Sleep disturbance
\end{itemize}
\textbf{Important:} If this symptom persists, your doctor will check for serious underlying causes.

\section*{Self-Care / Home Management}
\begin{itemize}
  \item Rest and avoid activities that make it worse.
  \item Maintain adequate hydration and a balanced diet.
  \item Over-the-counter remedies may provide symptomatic relief (consult a pharmacist or doctor).
  \item Monitor symptoms and seek professional advice if they persist or worsen.
  \item Avoid self-medicating with prescription medications without medical guidance.
\end{itemize}

\section*{How Your Doctor Will Evaluate You}
\begin{itemize}
  \item A thorough history and physical examination are the first steps in evaluation.
  \item Laboratory tests (blood work, urinalysis) may be ordered based on clinical suspicion.
  \item Imaging studies (X-ray, ultrasound, CT, MRI) may be indicated when structural pathology is suspected.
  \item Specialist referral is arranged when the diagnosis remains uncertain or requires specific expertise.
\end{itemize}

\section*{Treatment Options}
\begin{itemize}
  \item Treatment targets the underlying cause identified through diagnostic evaluation.
  \item Supportive care and symptom management are provided while awaiting definitive therapy.
  \item Medications, lifestyle modifications, and physical therapies may be prescribed as appropriate.
  \item Follow-up ensures treatment effectiveness and allows timely adjustments to the care plan.
\end{itemize}

\clearpage